%! Inverse Functions — Unit 1, Lesson 1.6 — Student Packet Cover Sheet
\documentclass[10pt]{article}
\usepackage{precalculus-article}
\usepackage{precalculus-boxes}
\usepackage{ltablex}
\keepXColumns

\begin{document}

% Full-bleed plum banner
\begin{tikzpicture}[remember picture, overlay]
  \fill[plum] (current page.north west)
    rectangle ([yshift=-0.9in]current page.north east);
\end{tikzpicture}
\vspace{-0.8in}

\begin{center}
  {\color{white}\textbf{\LARGE Precalculus}}\\[4pt]
  {\color{lilacmid}\large\textbf{Unit 1: Functions and Change}}\\[6pt]
  {\color{white}\textbf{\large Lesson 1.6 \quad Inverse Functions}}
\end{center}

\vspace{0.22in}
\namedateperiod
\vspace{0.18in}

\begin{learningtargetbox}
\small
\begin{enumerate}[leftmargin=1.5em, itemsep=5pt]
  \item I can reverse the input--output pairs of a function to build its
        inverse, and explain why the domain and range trade places.
  \item I can find $f^{-1}(x)$ by writing $y = f(x)$, swapping $x$ and $y$,
        and solving for $y$ --- and I know that $f^{-1}(x)$ is \textit{not}
        $\dfrac{1}{f(x)}$.
  \item I can verify a pair of inverses by composition, and restrict a domain
        when a function is not one-to-one.
\end{enumerate}
\end{learningtargetbox}

\vspace{0.14in}

\begin{tocbox}
\small
\renewcommand{\arraystretch}{1.5}
\begin{tabularx}{\linewidth}{@{} c l X r @{}}
\rowcolor{lilac}
\textbf{\#} & \textbf{Component} & \textbf{Description} & \textbf{Score} \\
\midrule
1 & Warm-Up        & Spiral review of prerequisite skills                & \blank{1.2cm} \\
2 & Guided Notes   & Building, checking, and reading inverse functions   & \blank{1.2cm} \\
3 & Group Activity & Reading the thermometer backwards                   & \blank{1.2cm} \\
4 & Exit Ticket    & Independent check on finding and verifying $f^{-1}$ & \blank{1.2cm} \\
5 & Homework       & Practice and extension                              & \blank{1.2cm} \\
\midrule
  & \multicolumn{2}{r}{\textbf{Total}} & \blank{1.2cm} \\
\end{tabularx}
\end{tocbox}

\vspace{0.14in}

\begin{remindbox}
\small
Lesson 1.5 fed one function into another. Today we build the function that
\textbf{undoes} another one --- and we use composition to prove it works: run a
function and its inverse back to back, in either order, and you land exactly
where you started. One warning carries the whole lesson: the $-1$ in $f^{-1}$
is a \textbf{name}, not an exponent.
\end{remindbox}

\end{document}
